% Part: incompleteness
% Chapter: incompleteness-provability
% Section: fixed-point-lemma

\documentclass[../../../include/open-logic-section]{subfiles}

\begin{document}

\olfileid{inc}{inp}{fix}

\olsection{لم نقطهٔ ثابت}

\begin{explain}
لم نقطهٔ ثابت می‌گوید برای هر !!{formula}~$!B(x)$، !!a{sentence}ی
چون~$!A$ وجود دارد که
$\Th{T} \Proves !A \liff !B(\gn{!A})$، مشروط بر اینکه~$\Th{T}$
گسترشی از~$\Th{Q}$ باشد. در مورد جملهٔ دروغ‌گو، می‌خواهیم~$!A$ (به‌طور
اثبات‌پذیر در~$\Th{T}$) با «$\gn{!A}$ کاذب است» هم‌ارز باشد؛ یعنی با
این گزاره که~$\Gn{!A}$ عدد گودل یک !!{sentence} کاذب است. برای فهم
ایدهٔ اثبات، مقایسهٔ آن با توضیح غیرصوری کواین دربارهٔ~$!A$ سودمند
است: «`اگر نقل‌قول خودش پیش از آن بیاید کذب می‌دهد' اگر نقل‌قول خودش
پیش از آن بیاید کذب می‌دهد.» می‌توان عملِ گرفتن یک عبارت و سپس
ساختن جمله‌ای با قراردادن نقل‌قول همان عبارت پیش از آن را
\emph{قطری‌کردن} عبارت، و نتیجه را قطری‌سازی آن نامید. پس قطری‌سازیِ
«اگر نقل‌قول خودش پیش از آن بیاید کذب می‌دهد» چنین است: «`اگر نقل‌قول
خودش پیش از آن بیاید کذب می‌دهد' اگر نقل‌قول خودش پیش از آن بیاید کذب
می‌دهد.» اکنون توجه کنید که جملهٔ دروغ‌گوی کواین قطری‌سازیِ «کذب
می‌دهد» نیست، بلکه قطری‌سازیِ «اگر نقل‌قول خودش پیش از آن بیاید کذب
می‌دهد» است. بنابراین خاصیتی که قطری می‌شود تا خودِ جملهٔ دروغ‌گو را
به دست دهد، خود دربردارندهٔ قطری‌سازی است!{}

در زبان حساب، نقل‌قولِ !!a{formula}ی دارای یک متغیر آزاد را با محاسبهٔ
عدد گودل آن و سپس جانشانی عددنمای استانداردِ آن عدد گودل به جای متغیر
آزاد می‌سازیم. قطری‌سازیِ~$!E(x)$ برابر~$!E(\num{n})$ است، که در آن
$n = \Gn{!E(x)}$. (از این پس، $\num{\Gn{!E(x)}}$ را به اختصار
$\gn{!E(x)}$ می‌نویسیم.) پس اگر~$!B(x)$ به معنای «کاذب است» باشد،
«اگر نقل‌قول خودش پیش از آن بیاید کذب می‌دهد» به معنای «وقتی بر عدد
گودل قطری‌سازیِ خودش اعمال شود کذب می‌دهد» خواهد بود. اگر نمادی چون
$\Obj{diag}$ برای تابع~$\fn{diag}(n)$ داشتیم که عدد گودل قطری‌سازیِ
!!{formula}ی با عدد گودل~$n$ را محاسبه می‌کند، می‌توانستیم~$!E(x)$ را
به صورت~$!B(\Obj{diag}(x))$ بنویسیم. آنگاه صورت کواینیِ جملهٔ دروغ‌گو
قطری‌سازیِ همین عبارت بود؛ یعنی~$!E(\gn{!E(x)})$ یا
$!B(\Obj{diag}(\gn{!B(\Obj{diag}(x))}))$. البته~$!B(x)$ اکنون
می‌تواند هر خاصیت دیگری باشد و همین ساخت کار خواهد کرد. برای قضیهٔ
ناتمامیت، $!B(x)$ را به معنای «$x$ در~$\Th{T}$ !!{derivable} نیست»
می‌گیریم. آنگاه~$!E(x)$ به معنای «وقتی بر عدد گودل قطری‌سازیِ خودش
اعمال شود، !!a{sentence}ی می‌دهد که در~$\Th{T}$ !!{derivable} نیست»
خواهد بود.

برای صوری‌کردن این مطلب در~$\Th{T}$، باید راهی برای صوری‌کردن
$\fn{diag}$ بیابیم. تابع~$\fn{diag}(n)$ محاسبه‌پذیر، و در واقع بازگشتی
اولیه است: اگر~$n$ عدد گودل فرمول~$!E(x)$ باشد، $\fn{diag}(n)$ عدد
گودل~$!E(\gn{!E(x)})$ را برمی‌گرداند. (به یاد آورید که~$\gn{!E(x)}$
عددنمای استانداردِ عدد گودل~$!E(x)$، یعنی~$\num{\Gn{!E(x)}}$، است.)
اگر~$\Obj{diag}$ نماد تابعی در~$\Th{T}$ و بازنمایندهٔ تابع
$\fn{diag}$ بود، می‌توانستیم~$!A$ را فرمول
$!B(\Obj{diag}(\gn{!B(\Obj{diag}(x))}))$ بگیریم. توجه کنید که
\begin{align*}
\fn{diag}(\Gn{!B(\Obj{diag}(x))}) & =
\Gn{!B(\Obj{diag}(\gn{!B(\Obj{diag}(x))}))} \\
& = \Gn{!A}.
\end{align*}
اگر فرض کنیم~$\Th{T}$ می‌تواند
\[
\Obj{diag}(\gn{!B(\Obj{diag}(x))}) = \gn{!A},
\]
را !!{derive} کند، آنگاه می‌تواند
$!B(\Obj{diag}(\gn{!B(\Obj{diag}(x))}))
\liff !B(\gn{!A})$ را نیز !!{derive} کند. اما سمت چپ بنا بر تعریف همان
$!A$ است.

البته~$\Obj{diag}$ عموماً نماد تابعی از~$\Th{T}$ نیست و بی‌گمان از
نمادهای~$\Th{Q}$ هم نیست. اما چون~$\fn{diag}$ محاسبه‌پذیر و
\emph{بازنمایی‌پذیر} است، در~$\Th{Q}$ بازنمایی آن با فرمولی چون
$!D_{\fn{diag}}(x,y)$ انجام می‌شود.
پس به جای نوشتن~$!B(\Obj{diag}(x))$ می‌توانیم
$\lexists[y][(!D_{\fn{diag}}(x,y) \land !B(y))]$ را بنویسیم. دیگر
بخش‌های اثباتِ طرح‌شده در بالا به همان صورت پیش می‌روند، و در واقع
همهٔ آن‌ها از همان~$\Th{Q}$ نتیجه می‌شوند.
\end{explain}

\begin{lem}
\ollabel{lem:fixed-point} فرض کنید~$!B(x)$ هر فرمولی با یک متغیر
آزاد~$x$ باشد. آنگاه جمله‌ای چون~$!A$ وجود دارد که
$\Th{Q} \Proves !A \liff !B(\gn{!A})$.
\end{lem}


\begin{proof}
با داشتن~$!B(x)$، بگذارید~$!E(x)$ فرمول
$\lexists[y][(!D_{\fn{diag}}(x,y) \land !B(y))]$ باشد و~$!A$ را
قطری‌سازی آن، یعنی فرمول~$!E(\gn{!E(x)})$، بگیرید.

چون~$!D_{\fn{diag}}$ تابع~$\fn{diag}$ را بازنمایی می‌کند و
$\fn{diag}(\Gn{!E(x)}) = \Gn{!A}$، نظریهٔ~$\Th{Q}$ می‌تواند عبارات زیر
را !!{derive} کند:
\begin{align}
  & !D_{\fn{diag}}(\gn{!E(x)}, \gn{!A}) \ollabel{repdiag1} \\
  & \lforall[y][(!D_{\fn{diag}}(\gn{!E(x)},y) \lif
  \eq[y][\gn{!A}])]. \ollabel{repdiag2}
\end{align}
اکنون نشان می‌دهیم~$\Th{Q} \Proves !A \liff !B(\gn{!A})$. به‌طور
غیرصوری استدلال می‌کنیم و تنها منطق و واقعیت‌های !!{derivable} در
$\Th{Q}$ را به کار می‌بریم.

نخست~$!A$، یعنی~$!E(\gn{!E(x)})$، را فرض کنید. با بازگشت به تعریف
$!E(x)$ می‌بینیم~$!E(\gn{!E(x)})$ همان عبارت زیر است:
\[
\lexists[y][(!D_{\fn{diag}}(\gn{!E(x)},y) \land !B(y))].
\]
چنین~$y$ای را در نظر بگیرید. چون
$!D_{\fn{diag}}(\gn{!E(x)},y)$، بنا بر~\olref{repdiag2} داریم
$y = \gn{!A}$. پس از~$!B(y)$، عبارت~$!B(\gn{!A})$ را به دست می‌آوریم.

اکنون~$!B(\gn{!A})$ را فرض کنید. بنا بر~\olref{repdiag1} داریم:
\begin{align*}
& !D_{\fn{diag}}(\gn{!E(x)}, \gn{!A}) \land !B(\gn{!A}).
\intertext{از آن نتیجه می‌شود که}
& \lexists[y][(!D_{\fn{diag}}(\gn{!E(x)},y) \land !B(y))].
\end{align*}
اما این همان~$!E(\gn{!E(x)})$، یعنی~$!A$، است.
\end{proof}

\begin{digress}
این اثبات را با اثبات لم نقطهٔ ثابت در نظریهٔ محاسبه‌پذیری مقایسه کنید.
تفاوت این است که اینجا می‌خواهیم \emph{گزاره‌ای} را بر حسب خودش تعریف
کنیم، در حالی که آنجا می‌خواستیم \emph{تابعی} را بر حسب خودش تعریف کنیم؛
گذشته از این تفاوت، ایده واقعاً یکی است.
\end{digress}

\begin{prob}
  !!^a{formula}~$!A(x)$ یک \emph{تعریف صدق} است هرگاه رابطهٔ
  $\Th{Q} \Proves !B \liff !A(\gn{!B})$ برای همهٔ
  !!{sentence}sی چون~$!B$ برقرار باشد. با استفاده از لم نقطهٔ ثابت
  نشان دهید هیچ !!{formula}ی تعریف صدق نیست.
\end{prob}

\end{document}
